\section{Compatible finite levels and the exact Siegel comparison}
\label{sec:pel-level-system}

The individual integral models can be coordinated without changing
their primitive line, coefficient field, or lattice.  The important
choice is to keep the rational cones fixed when the level changes.
Refining a fan would only give a proper transition, whereas the source
requires a finite one.

\begin{theorem}[The original simultaneous level clauses]
\label{thm:pel-level-system}
The original Siegel family clause \texttt{N0} holds, with $j=1$.
The individual models can be chosen simultaneously so that, for
each prime $p$ and each original valuation subring $W$ of $\mathbb C$
in which $p$ is a nonunit, the resulting local system satisfies
the complete clause \texttt{N2}.
In particular, every prescribed level and coefficient inclusion has
the exact finite \texttt{IsTransition}, and one auxiliary odd prime
works for every original level in the full \texttt{SiegelCompare}
quantifier.  These choices are compatible with the global models of
Theorem~\ref{thm:pel-global-domination}.
\end{theorem}

\subsection*{The precise Siegel moduli problem}

Let $g\geq1$, let $n$ be an odd prime different from $p$, and put
$R=\mathbb Z_{(p)}[\zeta_n]$, which is exactly the source's
$\texttt{ZP}(p,n)$.  Use the fine moduli scheme of principally
polarized abelian schemes of dimension $g$ with full symplectic
$n$-level and fixed Weil-pairing root $\zeta_n$.  It is a smooth
quasi-projective $R$-scheme of relative dimension $g(g+1)/2$
\cite[Theorem 1.4.1.11, Corollaries 1.4.1.12 and 7.2.3.10]{level-system-lan}.
The dimension is that of symmetric $g\times g$ deformation matrices.
Use the original theta-commutator convention
$e_\lambda(\alpha u,\alpha v)=\zeta_n^{\omega_{\rm std}(u,v)}$.
To identify it with Lan's opposite Hodge-sign convention, relabel
by the integral antisymplectic matrix $\operatorname{diag}(I_g,-I_g)$
and recover the original labels by its inverse.  This works for every
$g$, fixes $\zeta_n$, and changes neither the underlying lattice nor
the labelled points in the recovered object.

Write $(A,\lambda)$ for its universal polarized scheme and let
\[
                    P=(1\times\lambda)^*\mathcal P_A
                    \quad\hbox{on } A\times A.
\]
This is the source's line $P$, not a theta line on $A$.
Its diagonal pullback has polarization $2\lambda$.  Accordingly
its $n$-torsion commutator is $\zeta_n^{2\omega}$, precisely the
factor of two in \texttt{SiegPre.geom}.

We verify the source's apparently broader valuation-ring property.
A line $P$ rigidified on both axes of $A\times A$ determines, by
the rigidified Picard functor and the rigidity lemma, a homomorphism
$\lambda:A\to A^\vee$ with $P\simeq(1\times\lambda)^*\mathcal P_A$.
Normalize the two given axis trivializations to agree at $(0,0)$.
Symmetry of $P$ says $\lambda=\lambda^\vee$.
These statements hold over arbitrary bases by finite presentation
\cite[Sections 1.3.1--1.3.2]{level-system-lan}.
The source's trivial-line-kernel condition is exactly that
$\ker\lambda$ is the trivial group functor.  For a polarization,
this is equivalent to principality.  Conversely, a principal
polarization satisfies that condition on every test scheme.

Now let $V$ be any unit ball specified in \texttt{SiegelFamily.points},
with fraction field $C$.  The generic principal condition gives
an isomorphism $\lambda_C$.  Its inverse extends over $V$ by the
normal finite-approximation argument of the valuation-extension
theorem: descend the two abelian schemes, $\lambda$, its generic
inverse and their identities to a finite-type subalgebra of $V$
over $R$; normalize it inside its fraction field; and apply
normal-base full faithfulness for semiabelian schemes
\cite[Proposition 3.3.1.5]{level-system-lan}.
The normalization is finite because the original arithmetic ring
is excellent, and remains inside the integrally closed ring $V$.
Pullback gives the inverse over this very $V$.

The given ampleness on the closed fibre implies relative ampleness
of the diagonal line over the local base.  One may check this after
finite presentation: the ample-fibre locus is open for a proper
finitely presented family and contains the closed point, hence all
of the local spectrum.  Lan's Proposition 1.3.2.15 now says that
$\lambda$ is a polarization, including residue characteristic two.
The full $n$-division conditions in \texttt{SiegPre} give an
isomorphism of finite etale group schemes and force dimension $g$.
Since $n$ is odd, the identity for the commutator of $2\lambda$
uniquely recovers that for $\lambda$.  Thus the given object is an
object of the fine Siegel moduli problem over $V$, and is represented
by the required point.  This proves the complete \texttt{points}
field without a discreteness assumption.  The source's
dual-number \texttt{unram} field is immediate from the same fine
representing property.  This proves \texttt{N0} with $j=1$.
For $g=0$, when needed in \texttt{N2}, use $\operatorname{Spec}R$
and the trivial abelian scheme and line.

\subsection*{One rational fan for every level}

Choose an integral symplectic basis of the original unimodular
alternating lattice $(\Lambda,B)$.  Such a basis exists by successively
splitting a pair with pairing one and taking its orthogonal complement.
The basis identifies $B$ with the source's $\omega_{\rm std}$;
there is no rational change of lattice.

In the Siegel group on this fixed lattice, choose a smooth projective
rational cone system at principal level three, invariant under the
whole integral compact group $\operatorname{GSp}(\widehat\Lambda)$.
This is possible by
\cite[Proposition 2.8(1)--(2)]{level-system-projective}, since that
compact group normalizes principal level three.  We keep its rational
cones, rather than choosing new subdivisions at other levels.

Here smoothness at varying auxiliary Siegel levels presents no hidden
problem.  At a rational isotropic flag, the central unipotent matrices
have the form $\left(\begin{smallmatrix}1&B_0\\0&1\end{smallmatrix}\right)$
with $B_0$ symmetric.  Intersection with principal level $m$ replaces
the integral lattice of these matrices by $m$ times that lattice.
The dual cusp character lattice is therefore multiplied by $1/m$.
All the relevant lattices are homothetic.  A cone smooth for one
is smooth for each after rescaling its primitive generators.
The invariant system consequently defines smooth projective auxiliary
Siegel compactifications at every principal level $m\geq3$.
Projectivity is unchanged: multiply the rational support function
by a common denominator to make it integral at the chosen level.
Compatibility along faces and under the integral compact group
is unchanged as well.  This is the same-cone construction of
\cite[Construction 7.3.1.6 and the proof of Theorem 7.3.3.4]{level-system-lan}.

For an original level put $L=\operatorname{lcm}(M,A)$.
Over its original coefficient field, the half/full-level equivalence
identifies its group with the principal PEL subgroup $U(L)$.
At residue characteristic $p$, set $m=L/p^{v_p(L)}$.
The original \texttt{GoodAt} condition gives $m\geq3$.
Its image is contained in the auxiliary Siegel principal subgroup
of level $m$.  Use the singleton original symplectic lattice; the
auxiliary abelian scheme is the same abelian scheme, with its
principal polarization, not an isogenous replacement or a power.

Pull the fixed Siegel cones back to the PEL cones.  This induced
PEL system may be nonsmooth; that is allowed.  Define the integral
PEL compactification by normalization of the auxiliary Siegel
compactification in the original generic moduli function fields,
with all original generic components and the prescribed cyclotomic
component over the original coefficient ring.  These are exactly
the models in
\cite[(7.7)--(7.10), Theorem 7.14]{level-system-ramified}.
Theorem 7.14 identifies the result canonically whenever another
auxiliary level induces the same rational cones.  The individual
model, line, integral theta and valuation-extension proofs apply
to these choices.  In particular their complete
\texttt{LevelStructureConditions} are already available.

\subsection*{Finite transitions over the exact coefficient rings}

Suppose $l\mid l_2$ and $F\subset F_2$ as in the source.  Then
$L\mid L_2$.  Use for both models the lower auxiliary Siegel
level $m=L/p^{v_p(L)}$, which contains both PEL groups, with the
fixed rational cones.  Work over the larger original valuation ring
$O_2=c_2.V\,W$.  Over an affine open of this common ambient scheme,
let $A_0$ be its ring, let $B_1$ be the base-changed lower finite
normalization ring, and let $B_2$ be the upper normalization ring.
The generic forgetting map sends $B_1$ into the total fraction ring
of $B_2$.  Its image is integral over $A_0$, hence lies in $B_2$.
Moreover a finite set of $A_0$-module generators for $B_2$ also
generates it as a module over $B_1$.  Thus the resulting map
\[
                       Y_2\longrightarrow Y\times_O O_2
\]
is finite.  This argument does not require the base-changed lower
model to be normal.  Nor does it assert that the map of localized
number rings is finite.  It applies componentwise if necessary.

The maps glue uniquely by their generic restrictions.  Their
compositions agree for the same reason.  The tautological family
comparisons are cartesian on the common normalization presentation;
normal-base full faithfulness extends the action and the original
rigidified cubical-line comparisons.  On the generic moduli interior,
forgetting labels gives exactly
\[
             \eta_2(a)^{M_2/M}=\eta(a),\qquad
             \eta_{A,2}(a)^{A_2/A}=\eta_A(a).
\]
The half-level ideals cause no ambiguity: their stipulated inclusion
is precisely the inclusion that makes this forgetting map defined.
This supplies every field of \texttt{IsTransition}.  Its finite
morphism is after the stipulated coefficient base change, as required.
No claim of etaleness along the toroidal boundary is made.

\subsection*{One auxiliary prime for every original level}

Put $K_0=\mathbb Q(\sqrt{-d})$.  Since $d$ is prime and
$d\equiv7\pmod {16}$, its ring of integers is
$\mathbb Z[w_0]$.  The source's \texttt{Level.fromK} says
\[
      \mathfrak b=\mathfrak a\mathcal O_E,
      \qquad \mathfrak a=\mathfrak b\cap\mathcal O_{K_0}.
\]
The extension $\mathcal O_E/\mathcal O_{K_0}$ is finite faithfully
flat.  Contracting the two original ideal identities gives
\[
          \mathfrak a\overline{\mathfrak a}=(M),
          \qquad \mathfrak a+\overline{\mathfrak a}=\mathcal O_{K_0}.
\]
An inert prime is fixed by conjugation, so these equalities imply
that it does not divide $M$, for every original level.
This conclusion uses \texttt{fromK}; arbitrary half-ideals in
$E$ would not suffice.

Choose an odd inert prime $n\ne p$.  Such a prime can be chosen
outside any finite set without an analytic hypothesis.  For the
nontrivial quadratic character of $K_0$, choose an integer with
character $-1$, and use the Chinese remainder theorem to keep that
character while making the integer one modulo every excluded
unramified prime (including two and $p$ when applicable).
A prime factor with character $-1$ is then new.  Thus
$n$ is coprime to every $M$.  Take the Siegel family just constructed
at level $n$, so $j=1$, independently of the original level and
coefficient field.

\subsection*{The finite etale comparison on the abelian loci}

Let $l_1$ and $F_1$ be precisely those in \texttt{N2}:
$M_1=M$, $\mathfrak b_1=\mathfrak b$,
$A_1=\operatorname{lcm}(A,n)$, and $F_1=F(\zeta_n)$.
Let $U$ and $U_1$ be the actual abelian loci.
The extension of the original local number rings at $p$ is
essentially etale, since $n\ne p$.  Thus $U\times_OO_1$ is normal.
The missing compatible $n$-labels on its actual abelian scheme
form a finite etale cover.  The conditions of being symplectic,
$\mathcal O_E$-linear and of the prescribed liftable lattice type
select the original generic components of this cover; they are
open and closed.  Equivalently, take exactly the components whose
generic fibre is the original generic forgetting cover.
That generic cover is surjective, so the selected finite etale
cover is surjective over each component of the base.

The finite transition constructed above pulls back the original
family, so the inverse image of $U$ is exactly $U_1$.
Both $U_1$ and the selected cover are normal finite schemes over
$U\times_OO_1$ with the same generic function fields.  Uniqueness
of normalization identifies them.  This proves the complete finite
etale surjectivity assertion in \texttt{SiegelCompare}, including
the case $n\mid A$, when no new level is added.

Using the fixed integral symplectic basis, the exact full-$A_1$
label map gives the Siegel $n$-level
$a\mapsto\eta_{A,1}(a)^{A_1/n}$.  Its actual abelian scheme and
the line $\delta\mathcal L$ give the source's \texttt{SiegPre}.
The map
\[
                     U_1\longrightarrow \mathcal A_{g,n}\times_R O_1
\]
is finite: choose this very Siegel space as the auxiliary interior
in the normalization construction, as permitted by Theorem 7.14.
The genus is exactly $[E:\mathbb Q](p_D+q_D)/2$.
The product comparison preserves the biextension line and all
label powers, proving the specified \texttt{SiegPB} and square.

\subsection*{The reduced flat unramified intermediate model}

Assume now the extra source hypothesis $p\nmid\operatorname{lcm}(M,A)$.
Then every original label order, including $A_1$, is prime to $p$.
Over the fine Siegel scheme just used, consider the scheme of
$\mathcal O_E$-endomorphism structures with the required Rosati
involution, together with both entire original finite label maps and their
relations.  Endomorphism structures form a finite scheme by
\cite[Proposition 1.3.3.7]{level-system-lan}, and are formally
unramified by the rigidity of homomorphisms of abelian schemes
\cite[Lemma 2.2.2.1]{level-system-lan}.  The prime-to-$p$ label
schemes are finite etale, and imposing their identities is a closed
condition.  The resulting scheme $P$ is finite and formally
unramified over the Siegel base.

The original generic moduli space maps to $P$ as a monomorphism:
the abelian scheme, action and finite labels determine its object,
and the primitive line is recovered uniquely from $w_0$.
This morphism is finite, because the generic Siegel map is finite
and $P$ is separated over that base.  Thus it is a closed immersion.
Let $X'$ be its schematic closure in $P$ with the reduced structure.
Its generic fibre is the unchanged original generic space.  Every
component dominates $O_1$, so $X'$ is flat over this DVR.  It is
reduced, finite and formally unramified over the Siegel base.

The normalization of $X'$ in its generic fields is exactly $U_1$.
It is finite and surjective, and is an isomorphism over a dense
open $W\subset X'$; the inverse image of $W$ is dense as well.
The actual Siegel abelian scheme pulls back to $X'$.  Its universal
action, the primitive line $(1,\lambda\iota(w_0))^*\mathcal P$,
and both finite label maps give the source's \texttt{SAB} and
the two required extended sections.  Pulling back to $U_1$ gives
the original family and line, by their generic equality and
normal-base full faithfulness.  All the quantified generic-section
equalities follow from these equalities of morphisms.
These are precisely the remaining fields of \texttt{SiegelCompare}.

The construction proves the entire original \texttt{N2}, not only
proper level transitions or a generic Siegel map.  The remaining
formal-arc and numerical height, section and global norm obligations
are not supplied by this theorem.
